% ============================================================
% Manual Tecnico - Sistema de Liquidacion Tarifaria EPC
% Compilable a PDF con Overleaf o pdflatex
% Autor: Ikhthys Felipe Bernal Pachon
% Fecha: Agosto de 2026
% ============================================================

\documentclass[11pt,a4paper,oneside]{report}

% Paquetes basicos
\usepackage[utf8]{inputenc}
\usepackage[T1]{fontenc}
\usepackage[spanish]{babel}
\usepackage[margin=2.5cm]{geometry}
\usepackage{graphicx}
\usepackage{xcolor}
\usepackage{hyperref}
\usepackage{url}
\usepackage{booktabs}
\usepackage{longtable}
\usepackage{tabularx}
\usepackage{array}
\usepackage{listings}
\usepackage{textcomp}
\usepackage{amsmath}
\usepackage{amssymb}
\usepackage{enumitem}
\usepackage{fancyhdr}
\usepackage{titlesec}
\usepackage{multirow}
\usepackage{caption}
\usepackage{ragged2e}

% Column type para tablas: ragged-right con ancho fijo
% Previene que LaTeX hyphen palabra por palabra en columnas estrechas
\newcolumntype{L}[1]{>{\RaggedRight\arraybackslash}p{#1}}

% Configuracion de hyperref (sin colores chillones)
\hypersetup{
    colorlinks=true,
    linkcolor=blue!60!black,
    citecolor=blue!60!black,
    urlcolor=blue!60!black,
    pdftitle={Manual Tecnico - Sistema de Liquidacion Tarifaria EPC},
    pdfauthor={Ikhthys Felipe Bernal Pachon}
}

% Configuracion de listings (codigo)
\lstset{
    basicstyle=\ttfamily\small,
    keywordstyle=\color{blue!60!black}\bfseries,
    commentstyle=\color{green!50!black}\itshape,
    stringstyle=\color{red!60!black},
    showstringspaces=false,
    breaklines=true,
    frame=single,
    framerule=0.5pt,
    rulecolor=\color{black!30},
    backgroundcolor=\color{black!3},
    numbers=left,
    numberstyle=\tiny\color{black!50},
    numbersep=5pt
}

% Configuracion de titulos
\titleformat{\chapter}[hang]{\Huge\bfseries}{\thechapter}{1em}{}
\titleformat{\section}[hang]{\Large\bfseries}{\thesection}{1em}{}
\titleformat{\subsection}[hang]{\large\bfseries}{\thesubsection}{1em}{}

% Configuracion de headers/footers
\pagestyle{fancy}
\fancyhf{}
\fancyhead[L]{\small Manual Tecnico EPC}
\fancyhead[R]{\small Ikhthys Felipe Bernal Pachon}
\fancyfoot[C]{\thepage}
\renewcommand{\headrulewidth}{0.4pt}
\renewcommand{\footrulewidth}{0.4pt}

% Configuracion de tablas
\renewcommand{\arraystretch}{1.2}

% ============================================================
% Metadatos
% ============================================================

\title{
    \vspace{-2cm}
    \Huge \textbf{Sistema de Liquidacion Tarifaria EPC} \\[0.5cm]
    \LARGE \textbf{Manual Tecnico} \\[1cm]
    \normalsize Trabajo de grado --- Ingenieria de Sistemas \\[0.3cm]
    \normalsize Universidad de Cundinamarca \\[0.3cm]
    \normalsize Cliente piloto: Empresas Publicas de Cundinamarca (EPC)
}

\author{
    \Large \textbf{Ikhthys Felipe Bernal Pachon} \\[0.5cm]
    \normalsize Ingenieria de Sistemas \\
    \normalsize Universidad de Cundinamarca \\
    \normalsize ikhthys.bernal@example.edu \\ % Reemplazar por email real antes de compilar
}

\date{Agosto de 2026}

% ============================================================
% Comienzo del documento
% ============================================================

\begin{document}

\maketitle

\thispagestyle{empty}

% ============================================================
% Resumen
% ============================================================

\begin{abstract}
\noindent
Este manual describe la arquitectura, implementacion y estado del Sistema de Liquidacion Tarifaria para Empresas Publicas de Cundinamarca (EPC). El sistema automatiza la captura de lecturas de medidores en campo (modo offline) y el calculo de liquidaciones tarifarias segun la metodologia CRA vigente para prestadores rurales con menos de 5.000 suscriptores.

\noindent
El proyecto fue desarrollado como trabajo de grado en Ingenieria de Sistemas de la Universidad de Cundinamarca, con cliente piloto EPC. El alcance incluye: aplicacion mobile offline-first (Expo + React Native), backend API REST (.NET 8 + PostgreSQL), dashboard web (HTML estatico), y motor tarifario con cumplimiento de la Res CRA 825/2017, Res CRA 907/2019 y Res CRA 750/2016.

\noindent
Al cierre de esta entrega, el sistema cuenta con 2.123 tests automatizados verdes, 8 de 10 gaps de compliance cerrados y una arquitectura limpia basada en dominio puro compartido entre capas.
\end{abstract}

\tableofcontents
\listoftables
\listoffigures

\newpage

% ============================================================
% Capitulo 1: Introduccion
% ============================================================

\chapter{Introduccion}

\section{Proposito del documento}

Este manual tecnico esta dirigido al equipo de EPC que recibira el proyecto y a cualquier ingeniero que continue su desarrollo. Su proposito es servir como referencia completa de:

\begin{itemize}[leftmargin=*]
    \item La arquitectura del sistema y sus componentes.
    \item El estado actual del desarrollo y los gaps pendientes.
    \item Las decisiones arquitectonicas clave y su rationale.
    \item El cumplimiento regulatorio CRA 825/2017 y resoluciones complementarias.
    \item El flujo de desarrollo (SDD workflow + TDD-strict).
    \item Los issues conocidos y los proximos pasos sugeridos.
\end{itemize}

\section{Alcance del proyecto}

El Sistema de Liquidacion Tarifaria EPC automatiza:

\begin{enumerate}[leftmargin=*]
    \item \textbf{Captura en campo:} lectura de medidores, fotografia de evidencia y verificacion de estrato socioeconomico.
    \item \textbf{Calculo tarifario:} aplicacion de la metodologia CRA 825/2017 con factores de subsidio (estratos 1-3), contribuciones (estratos 5-6) y minimo vital.
    \item \textbf{Sincronizacion offline-first:} cola idempotente con orden de dependencias (suscriptor $\to$ medidor $\to$ lectura $\to$ liquidacion).
    \item \textbf{Consulta administrativa:} dashboard web para consulta de suscriptores, lecturas y liquidaciones.
\end{enumerate}

\section{Audiencia}

\begin{itemize}[leftmargin=*]
    \item \textbf{Equipo de EPC:} personal tecnico que operara y mantendra el sistema.
    \item \textbf{Ingenieros de mantenimiento:} continuaran el desarrollo y abordaran los gaps pendientes.
    \item \textbf{Auditoria CRA:} verificacion del cumplimiento normativo del sistema.
\end{itemize}

\section{Contexto institucional}

\begin{itemize}[leftmargin=*]
    \item \textbf{Cliente piloto:} Empresas Publicas de Cundinamarca (EPC).
    \item \textbf{Caso de uso:} 2-3 prestadores rurales durante 3-5 dias de prueba.
    \item \textbf{Tipo:} Trabajo de grado, Ingenieria de Sistemas, Universidad de Cundinamarca.
    \item \textbf{Ano:} 2026.
\end{itemize}

% ============================================================
% Capitulo 2: Vision general del sistema
% ============================================================

\chapter{Vision general del sistema}

\section{Descripcion general}

El sistema es una solucion full-stack para prestadores rurales de acueducto en Colombia. Su objetivo es reemplazar el calculo manual de liquidaciones tarifarias por un flujo digital que cumpla con la normativa CRA y permita la captura en campo sin depender de conectividad.

\section{Usuarios principales}

\begin{table}[h]
\centering
\caption{Tipos de usuario del sistema}
\begin{tabularx}{\textwidth}{L{3.5cm} L{4.5cm} X}
\toprule
\textbf{Tipo} & \textbf{Rol} & \textbf{Pantallas principales} \\
\midrule
Operario de campo & Captura lecturas y fotografias en dispositivos Android & RutaDeHoy, CapturarLectura, DetalleSuscriptor \\
Administrador & Configura parametros tarifarios y gestiona operarios & ParametrosTarifa, Admin Operarios \\
Supervisor & Consulta dashboard y valida datos & Dashboard web \\
\bottomrule
\end{tabularx}
\end{table}

\section{Flujo de operacion}

\begin{enumerate}[leftmargin=*]
    \item El operario de campo abre la app mobile (sin internet) y selecciona su prestador.
    \item Recorre la ruta del dia (lista de suscriptores asignados).
    \item Captura lectura + fotografia del medidor en cada suscriptor.
    \item La app calcula la liquidacion en tiempo real con la metodologia CRA.
    \item Cuando hay WiFi disponible, el operario dispara la sincronizacion manual.
    \item El backend .NET recibe los datos via API REST con protocolo idempotente.
    \item El dashboard web consulta los datos para supervision administrativa.
\end{enumerate}

\section{Stack tecnologico}

\begin{table}[h]
\centering
\caption{Stack tecnologico por capa}
\begin{tabularx}{\textwidth}{L{5cm} X}
\toprule
\textbf{Capa} & \textbf{Tecnologia} \\
\midrule
Lenguaje nucleo & TypeScript 5/6 strict \\
Workspaces & npm workspaces (monorepo) \\
Dominio puro & TypeScript + Jest 30 + ts-jest 29 \\
Mobile & React Native 0.81, Expo SDK 54, React 19.1, expo-router 6, expo-sqlite 16 \\
Estado mobile & Zustand 5 \\
UI mobile & react-native-paper 5 \\
Hardware mobile & expo-camera, react-native-ble-plx (Bluetooth para impresion) \\
Backend & .NET 8, ASP.NET Core Minimal API, EF Core 8 \\
Base de datos servidor & PostgreSQL 16 \\
Base de datos local mobile & SQLite (expo-sqlite) \\
Base de datos local backend & better-sqlite3 (Node.js, en tests) \\
Dashboard & HTML5 + Vanilla JS + bcrypt.js (en cliente) \\
Especificacion & OpenSpec + Engram (memoria persistente cross-session) \\
Testing mobile & jest-expo + @testing-library/react-native \\
\bottomrule
\end{tabularx}
\end{table}

% ============================================================
% Capitulo 3: Arquitectura
% ============================================================

\chapter{Arquitectura del sistema}

\section{Principio arquitectonico: separacion de capas estricta}

El sistema aplica \textbf{separacion de capas estricta} como principio fundamental:

\begin{enumerate}[leftmargin=*]
    \item \textbf{Dominio puro:} logica de negocio en TypeScript puro, sin dependencias de framework, plataforma ni base de datos.
    \item \textbf{Puertos (interfaces):} contratos que el dominio expone para acceso a datos.
    \item \textbf{Adaptadores:} implementaciones concretas de los puertos para cada plataforma.
    \item \textbf{UI (Pantallas):} unicamente presentacion; no contiene logica de negocio.
\end{enumerate}

\textbf{Implicancia practica:} si una regla de calculo cambia, se modifica una sola vez en el dominio y se propaga automaticamente. NO se duplica logica entre capas.

\section{Diagrama de componentes}

\begin{figure}[h]
\centering
\begin{tabular}{|c|}
\hline
\textbf{Dispositivo movil (Android)} \\
\hline
Aplicacion Expo / React Native \\
\hline
Pantallas (UI presentacional) \\
$\downarrow$ usa \\
Composicion (bootstrap, DI) \\
$\downarrow$ usa puertos del dominio \\
Dominio puro TypeScript compartido \\
$\downarrow$ usa interfaces de persistencia \\
Adaptadores expo-sqlite (repositorios) \\
$\downarrow$ usa \\
SQLite local \\
\hline
\end{tabular}
\quad
\begin{tabular}{|c|}
\hline
\textbf{Servidor (red LAN)} \\
\hline
MediApp.Api (.NET 8 Minimal API) \\
\hline
/api/v1/suscriptores \\
/api/v1/medidores \\
/api/v1/lecturas \\
/api/v1/liquidaciones \\
/api/v1/operarios \\
/health \\
wwwroot/ (Dashboard HTML estatico) \\
\hline
EF Core + Npgsql \\
$\downarrow$ \\
PostgreSQL 16 \\
\hline
\end{tabular}

\caption{Arquitectura de componentes del sistema}
\end{figure}

\section{Sincronizacion cliente-servidor}

La comunicacion entre mobile y backend es \textbf{manual} (no push automatico). El operario la dispara desde la pantalla Sync. El protocolo es idempotente:

\begin{enumerate}[leftmargin=*]
    \item El cliente envia cada item con un \texttt{id\_cliente} unico (formato \texttt{dispositivoId:idLocal}).
    \item El servidor valida y busca en la tabla \texttt{sync\_registros}.
    \item Si existe, retorna 200 OK sin modificar.
    \item Si no existe, inserta y registra, retorna 201 Created.
\end{enumerate}

Esto permite reintentar cualquier sincronizacion fallida sin riesgo de duplicados.

% ============================================================
% Capitulo 4: Componentes principales
% ============================================================

\chapter{Componentes principales}

\section{Dominio puro (TypeScript compartido)}

El dominio es TypeScript puro, sin dependencias de framework. Puede ejecutarse en React Native, Node.js o un navegador web. Se ubica en \texttt{mobile/src/\{modulo\}/}.

\begin{table}[h]
\centering
\caption{Modulos del dominio}
\begin{tabularx}{\textwidth}{L{4.5cm} X}
\toprule
\textbf{Modulo} & \textbf{Responsabilidad} \\
\midrule
\texttt{motor-tarifario/} & Calculo CRA: CF, CC, subsidios, validaciones \\
\texttt{parametros-tarifa/} & Dominio ParametrosTarifa con flag \texttt{aplica\_cmaa} e IPC editable \\
\texttt{acuerdo-municipal/} & AcuerdoMunicipal y estado del ciclo de vida \\
\texttt{suscriptores/} & Suscriptores y verificacion de estrato \\
\texttt{medidores/} & Medidores \\
\texttt{lecturas/} & Lecturas \\
\texttt{operarios/} & Operarios (admin y campo) \\
\texttt{prestadores/} & Prestadores (multi-tenant) \\
\texttt{factura/} & Factura y snapshot \texttt{validacion\_ambito} \\
\texttt{importacion/} & Importacion masiva desde CSV \\
\texttt{sincronizacion/} & Cola de sincronizacion offline \\
\texttt{captura-lecturas/} & Caso de uso: capturar lectura \\
\texttt{calculo/} & Calculos auxiliares \\
\texttt{auditoria/} & Trazabilidad regulatoria \\
\texttt{persistencia/} & Adaptadores SQLite (integration tests) \\
\texttt{cliente-http/} & HTTP client para backend .NET \\
\texttt{categorias-uso/} & Categorias (residencial, comercial, etc.) \\
\texttt{periodos/} & Periodos de facturacion \\
\texttt{componentes/} & Componentes auxiliares \\
\texttt{shared/} & Utilidades compartidas y ports \\
\bottomrule
\end{tabularx}
\end{table}

\section{Aplicacion mobile (Expo + React Native)}

La aplicacion mobile esta construida con Expo SDK 54 + React Native 0.81. Su estructura interna:

\begin{itemize}[leftmargin=*]
    \item \texttt{pantallas/}: pantallas de la aplicacion (ParametrosTarifa, RutaDeHoy, etc.).
    \item \texttt{persistencia/}: repositorios expo-sqlite.
    \item \texttt{composition/}: bootstrap + inyeccion de dependencias.
    \item \texttt{componentes/}: componentes UI presentacionales.
    \item \texttt{hooks/}: hooks custom (useParametrosFormState, etc.).
    \item \texttt{theme/}: tema y design tokens.
    \item \texttt{adapters/}: adapters (Bluetooth, impresion).
    \item \texttt{navegacion/}: navegacion + auth gate.
\end{itemize}

\paragraph{Pantallas principales:}
RutaDeHoy, CapturarLectura, DetalleSuscriptor, ParametrosTarifa, Sincronizacion, Configuracion, SetupInicial, Login.

\section{Backend (.NET 8)}

El backend es una Minimal API en .NET 8 con Entity Framework Core 8 + Npgsql. Su estructura:

\begin{itemize}[leftmargin=*]
    \item \texttt{Features/}: organizacion por bounded context (Operarios, Suscriptores, Medidores, Lecturas, Liquidaciones).
    \item \texttt{Persistence/}: DbContext de EF Core + migraciones PostgreSQL.
    \item \texttt{wwwroot/}: dashboard web HTML estatico.
\end{itemize}

\paragraph{Endpoints principales:}
\texttt{GET /health}, \texttt{POST /api/v1/suscriptores}, \texttt{GET /api/v1/suscriptores}, \texttt{POST /api/v1/medidores}, \texttt{GET /api/v1/medidores}, \texttt{POST /api/v1/lecturas}, \texttt{GET /api/v1/lecturas}, \texttt{POST /api/v1/liquidaciones}, \texttt{GET /api/v1/liquidaciones}, \texttt{POST /api/v1/operarios}, \texttt{GET /api/v1/operarios}, \texttt{PUT /api/v1/operarios/\{id\}}, \texttt{PATCH /api/v1/operarios/vincular-dispositivo}.

\section{Dashboard web}

El dashboard es una SPA estatica servida directamente por el backend. Caracteristicas:

\begin{itemize}[leftmargin=*]
    \item HTML5 + CSS3 + Vanilla JavaScript (sin frameworks).
    \item bcrypt.js cargado desde CDN para hash de contrasenas en cliente.
    \item Fetch API para comunicacion con el backend REST.
\end{itemize}

Secciones: Operarios (CRUD + vincular dispositivo), Suscriptores (consulta), Lecturas (consulta con datos del medidor y suscriptor), Liquidaciones (consulta).

% ============================================================
% Capitulo 5: Motor tarifario CRA
% ============================================================

\chapter{Motor tarifario CRA}

\section{Normativa aplicada}

\begin{table}[h]
\centering
\caption{Normativa CRA aplicada al sistema}
\begin{tabularx}{\textwidth}{L{3cm} L{3.5cm} X}
\toprule
\textbf{Norma} & \textbf{Aplica a} & \textbf{Alcance en el sistema} \\
\midrule
Res CRA 825/2017 & Prestadores rurales con menos de 5.000 suscriptores & Base normativa: formula CF, CC, subsidios E1-E3, contribuciones E5-E6, CMOG minimo \\
Res CRA 907/2019 & Modifica los articulos 9 y 10 de la 825 & Introduce CMAA (inversiones ambientales) y CMVIAA \\
Res CRA 750/2016 & Consumo basico por altitud & Limite basico varia con altitud\_msnm del prestador \\
Ley 142/1994 art. 99.6 & Topes de subsidios y contribuciones & Estrato 1 [-60\%], 2 [-50\%], 3 [-40\%], 5 [+50\%], 6 [+60\%] \\
Res CRA 881/2019 & Inversiones ambientales & Adiciona articulos sobre inversiones \\
\bottomrule
\end{tabularx}
\end{table}

\textbf{No aplica} a este proyecto: Res CRA 1032/2026 (vigente desde 24/03/2026) porque es para prestadores >5000 suscriptores urbanos.

\section{Funcion principal}

\begin{lstlisting}[language=Java,caption={Interfaz principal del motor},label={lst:motor}]
calcularLiquidacion(entrada: EntradaCalculo): ResultadoCalculo
\end{lstlisting}

\section{Parametros de entrada}

\begin{table}[h]
\centering
\caption{Campos de EntradaCalculo}
\begin{tabularx}{\textwidth}{L{4.5cm} L{2.5cm} X}
\toprule
\textbf{Campo} & \textbf{Tipo} & \textbf{Descripcion} \\
\midrule
\texttt{lecturaAnterior} & number & Lectura del periodo anterior, en metros cubicos \\
\texttt{lecturaActual} & number & Lectura del periodo actual, en metros cubicos \\
\texttt{estrato} & Estrato (1-6) & Estrato socioeconomico del suscriptor \\
\texttt{aplicaSubsidio} & boolean & Si aplica el factor de subsidio o contribucion \\
\texttt{periodo} & \{mes, anio\} & Periodo de facturacion \\
\texttt{parametros.cargoFijo} & number & CF mensual en pesos (puede ser 0 por minimo vital) \\
\texttt{parametros.precioM3} & number & Precio del metro cubico en bloque basico \\
\texttt{parametros.precioM3Excedente} & number & Precio del metro cubico en bloque excedente \\
\texttt{parametros.consumoBasico} & number & Limite del bloque basico (default: 20 metros cubicos) \\
\texttt{parametros.aplicaCmaa} & boolean & Opt-in para CMAA (inversiones ambientales) \\
\bottomrule
\end{tabularx}
\end{table}

\section{Formula de calculo}

\begin{lstlisting}[caption={Formula del motor tarifario CRA},label={lst:formula}]
consumo = lecturaActual - lecturaAnterior
consumoBasico = min(consumo, limiteBasico)
consumoExcedente = max(consumo - limiteBasico, 0)

cargoConsumo = round(consumoBasico * precioM3)
cargoExcedente = round(consumoExcedente * precioM3Excedente)
cargoFijo = round(cargoFijoRaw)

// Factor de estrato (CRA 825/2017):
// E1: -60% | E2: -50% | E3: -40% | E4: 0% | E5: +50% | E6: +60%

// Subsidio (E1-E3): sobre CF + cargoConsumo
subsidio = round(|factor| * (cargoFijo + cargoConsumo))

// Contribucion (E5-E6): sobre CF + consumo total
contribucion = round(factor * (cargoFijo + cargoConsumo + cargoExcedente))

// CMAA (opt-in): se agrega al total si aplicaCmaa = true
total = cargoFijo + cargoConsumo + cargoExcedente
        - subsidio + contribucion + cmAA
\end{lstlisting}

\section{Validaciones del motor}

\begin{itemize}[leftmargin=*]
    \item Lecturas no pueden ser negativas.
    \item Lectura actual $\geq$ lectura anterior.
    \item Cargo fijo $\geq$ 0 (0 es valido por minimo vital).
    \item Precio m\textsuperscript{3} y precio m\textsuperscript{3} excedente $> 0$.
    \item Limite consumo basico $> 0$.
    \item Estrato entre 1 y 6 (si se provee).
    \item Mes entre 1 y 12, anio $\geq 2000$.
\end{itemize}

% ============================================================
% Capitulo 6: Compliance CRA 825/2017
% ============================================================

\chapter{Compliance CRA 825/2017}

\section{Estado al cierre de esta entrega}

Resultado del gap analysis + 3 fases de implementacion:

\begin{table}[h]
\centering
\caption{Fases del compliance CRA 825}
\begin{tabularx}{\textwidth}{L{2.5cm} L{4.5cm} X c}
\toprule
\textbf{Fase} & \textbf{Change} & \textbf{Gaps cerrados} & \textbf{Estado} \\
\midrule
1 (julio 2026) & param-tarifa-res-825-compliance-phase1 & CMA minimo, IPC, anio base 2016, campo aps & Archivado \\
2 (agosto 2026) & param-tarifa-res-825-compliance-phase2 & Formula CF corregida, CMAA, validarAmbito, validarCmogMinimo, Acuerdo.estado, estado\_verificacion, 3 campos docs, @deprecated calcularCCUnitario & Archivado \\
3 (agosto 2026) & param-tarifa-residuales-cra-825 & Flag \texttt{aplica\_cmaa} (opt-in explicito), inputs IPC editables, unificacion minimo vital Opcion A, migracion 030 aditiva idempotente, refactors a modulos puros & Archivado \\
\bottomrule
\end{tabularx}
\end{table}

\section{Gaps cerrados (8)}

\begin{enumerate}[leftmargin=*]
    \item CMA minimo (Costo Medio de Administracion).
    \item IPC editable con anio base 2016.
    \item Campo \texttt{aps} (perdidas tecnicas) en el modelo.
    \item Formula CF corregida (Cargo Fijo recalculado).
    \item CMAA (inversiones ambientales) con flag \texttt{aplica\_cmaa} opt-in.
    \item Validadores \texttt{validarAmbito} y \texttt{validarCmogMinimo}.
    \item Minimo vital unificado (Opcion A: cargo fijo puede ser 0).
    \item Migracion aditiva idempotente para flag \texttt{aplica\_cmaa}.
\end{enumerate}

\section{Gaps diferidos explicitamente (2)}

\begin{table}[h]
\centering
\caption{Gaps diferidos}
\begin{tabularx}{\textwidth}{L{4.5cm} L{3.5cm} X c}
\toprule
\textbf{Gap} & \textbf{Cambio futuro} & \textbf{Razon del diferimiento} \\
\midrule
GAP-6: MetadataCalculo completo §11 (20+ campos) & auditoria-mejorada-cra-825 & Requiere 20+ campos adicionales de trazabilidad. Cambio dedicado por scope. \\
GAP-10: Contrato de salida del agente IA §12 & agente-validacion-cra-825 & Depende de decisiones de producto sobre integracion de IA. \\
\bottomrule
\end{tabularx}
\end{table}

\section{Porcentaje de cumplimiento}

\textbf{80\%} (8 de 10 gaps cerrados). NO dar por cerrado el compliance hasta abordar GAP-6 y GAP-10.

% ============================================================
% Capitulo 7: Decisiones arquitectonicas
% ============================================================

\chapter{Decisiones arquitectonicas clave}

\section{Dominio compartido entre mobile y otros clientes}

\textbf{Decision:} todo el dominio vive dentro de \texttt{mobile/src/\{modulo\}/}. NO se duplica logica en backend .NET.

\textbf{Rationale:} evita drift entre lo que calcula el mobile y lo que muestra el dashboard. Si la regla cambia, cambia en un solo lugar.

\textbf{Tradeoff:} el backend .NET NO tiene logica de calculo propia, solo persiste resultados que el mobile le envia.

\section{Backend NO hashea contrasenas}

\textbf{Decision:} el cliente (mobile con SHA-256 o dashboard con bcrypt) hashea antes del POST. El backend solo valida.

\textbf{Rationale:} simplifica el backend y elimina un vector de ataque (el backend nunca ve la contrasena en claro). El cliente ya tiene el hash necesario para login offline.

\textbf{Tradeoff:} el hash esta en el cliente, lo cual es un riesgo si el cliente es comprometido. Por eso el mobile usa SHA-256 y el dashboard bcrypt.

\section{Mobile offline-first con sync manual}

\textbf{Decision:} la captura funciona sin internet. La sincronizacion contra el backend es \textbf{manual} (el operario la dispara).

\textbf{Rationale:} los prestadores rurales tienen conectividad intermitente. Asumir conexion online es poco realista.

\textbf{Tradeoff:} requiere cola de sincronizacion con orden de dependencias y protocolo idempotente.

\section{Multi-tenancy por \texttt{id\_prestador}}

\textbf{Decision:} cada operario entra a SU prestador. Todas las tablas tienen FK \texttt{id\_prestador}.

\textbf{Rationale:} permite que un mismo deployment sirva a multiples prestadores sin filtrar datos.

\textbf{Tradeoff:} todas las queries deben incluir \texttt{WHERE id\_prestador = ?} (es una convencion disciplinada, no forzada por DB).

\section{Tests colocados, no centralizados}

\textbf{Decision:} cada modulo tiene su \texttt{\_\_tests\_\_/} adyacente, no hay un \texttt{\_\_tests\_\_/} raiz centralizado.

\textbf{Rationale:} facilita encontrar los tests junto al codigo. Reduce la friccion de agregar tests.

\textbf{Tradeoff:} requiere disciplina para no duplicar setup.

\section{Strict TDD mode}

\textbf{Decision:} cada test rojo se commitea antes del feature. Cada RED commit es atomico.

\textbf{Rationale:} historial de git cuenta la historia del desarrollo. El reviewer puede ver exactamente que se probo.

\textbf{Tradeoff:} mas commits por feature (tipicamente 2x: RED + GREEN). El ruido se mitiga con squash al mergear.

\section{Conventional commits en ingles}

\textbf{Decision:} mensajes de commit en ingles, formato conventional commits, sin \texttt{Co-Authored-By}, sin emojis.

\textbf{Rationale:} permite generar changelogs automaticamente y mantener historial buscable.

\textbf{Tradeoff:} requiere disciplina. El subject $\leq$ 50 caracteres fuerza a ser conciso.

% ============================================================
% Capitulo 8: Estado actual y metricas
% ============================================================

\chapter{Estado actual y metricas}

\section{Estado por componente}

\begin{table}[h]
\centering
\caption{Estado por componente del sistema}
\begin{tabularx}{\textwidth}{L{5cm} X c}
\toprule
\textbf{Componente} & \textbf{Estado} & \textbf{Verde?} \\
\midrule
Dominio (TypeScript compartido) & Completo, compliance CRA 825 cerrado & Si \\
Mobile (Expo + React Native) & MVP funcional con compliance cerrado & Parcial \\
Backend (.NET 8) & MVP funcional, sin deploy productivo & Parcial \\
Dashboard web & MVP funcional & Si \\
Tests automatizados & 2123 verde en 171 suites & Si \\
Type check & \texttt{tsc --noEmit} sin errores & Si \\
\bottomrule
\end{tabularx}
\end{table}

\section{Metricas finales}

\begin{table}[h]
\centering
\caption{Metricas finales del proyecto}
\begin{tabularx}{\textwidth}{L{5cm} X}
\toprule
\textbf{Metrica} & \textbf{Valor} \\
\midrule
Tests automatizados & 2123 verde (171 suites) \\
Archivos del dominio & repartidos en \texttt{mobile/src/} (modulos por dominio) \\
Pantallas mobile & \texttt{mobile/src/pantallas/} \\
Persistencia local & SQLite (\texttt{expo-sqlite} y \texttt{better-sqlite3} en tests) \\
Persistencia servidor & PostgreSQL 16 \\
Sync & manual desde pantalla Sync, protocolo idempotente con \texttt{sync\_registros} \\
Branches & \texttt{main}, \texttt{desarrollo} y \texttt{merge-desarrollo} (cleanup pendiente) \\
SDDs archivados & 22+ en \texttt{openspec/changes/archive/} \\
\bottomrule
\end{tabularx}
\end{table}

% ============================================================
% Capitulo 9: Issues conocidos
% ============================================================

\chapter{Issues conocidos}

\section{Criticos (bloquean produccion)}

Ninguno al cierre de esta entrega.

\section{Importantes (deben abordarse pronto)}

\begin{enumerate}[leftmargin=*]
    \item \textbf{Token de sesion fake-token.} El token actual es \texttt{fake-token-\{timestamp\}} con expiracion hardcoded a 24 horas. NO usar como mecanismo de seguridad en produccion. Migracion a GUID generado por backend .NET con \texttt{expiresAt} validado contra PostgreSQL.
    \item \textbf{Backend sin deploy productivo.} Esta configurado para dev local. NO hay pipeline CI/CD ni configuracion de hosting. Esta decision queda para EPC.
\end{enumerate}

\section{Menores (limpieza)}

\begin{enumerate}[leftmargin=*]
    \item \textbf{Directorio \texttt{src/\_\_tests\_\_/} residual.} Existe en la raiz del proyecto como leftover del SDD \texttt{delete-legacy-src-mirror}. NO contiene tests activos, pero es ruido. Considerar agregar a \texttt{.gitignore} o borrar.
    \item \textbf{Branch \texttt{merge-desarrollo} pendiente de cerrar.} Ya esta mergeada a \texttt{main}. Limpieza: \texttt{git push origin --delete merge-desarrollo}.
    \item \textbf{SDDs activos sin cerrar.} Hay 3 SDDs abiertos en \texttt{openspec/changes/} (\texttt{eliminar-campo-calle}, \texttt{factura-compliance-fase1}, \texttt{factura-compliance-polish}). Decidir si cerrarlos o continuar trabajo.
\end{enumerate}

\section{Optimizaciones pendientes}

\subsection*{Performance}
\begin{itemize}[leftmargin=*]
    \item Bundle size de Expo. Evaluar tree-shaking y lazy loading de pantallas no usadas.
    \item Query N+1 en dashboard. El listado de lecturas con medidor y suscriptor podria optimizarse con JOINs eager.
    \item Sync incremental. Actualmente la sync envia todos los pendientes. Evaluar sync incremental con cursor.
\end{itemize}

\subsection*{Code quality}
\begin{itemize}[leftmargin=*]
    \item ParametrosTarifa.tsx sigue en 456 lineas (post-decompose). Podria continuar decomponiendo subcomponentes grandes si crecen.
    \item \texttt{composicion/} vs \texttt{composition/} --- naming consistente.
    \item Migraciones SQL inline (889 lineas). Considerar extraer a archivos \texttt{.sql} separados.
\end{itemize}

\subsection*{Testing}
\begin{itemize}[leftmargin=*]
    \item Cobertura de UI. Tests actuales cubren dominio + repos. UI esta sub-testada.
    \item E2E tests. No hay tests end-to-end automatizados. Considerar Detox o Maestro.
    \item Mutation testing. Considerar Stryker para validar calidad de los tests.
\end{itemize}

% ============================================================
% Capitulo 10: Proximos pasos
% ============================================================

\chapter{Proximos pasos sugeridos}

\section{No urgentes pero recomendados}

\begin{enumerate}[leftmargin=*]
    \item Cerrar la branch \texttt{merge-desarrollo} local y remote (cleanup puro, ya mergeada a main).
    \item Abordar GAP-6 (\texttt{auditoria-mejorada-cra-825}) y GAP-10 (\texttt{agente-validacion-cra-825}) cuando el equipo lo priorice.
    \item Migracion del token fake-token a GUID real con expiracion validada contra PostgreSQL.
    \item Pipeline CI/CD productivo para backend .NET (actualmente solo mobile esta automatizado).
    \item Evaluar el SDD activo \texttt{factura-compliance-fase1} y \texttt{factura-compliance-polish} en \texttt{openspec/changes/}.
    \item Tests E2E con Detox o Maestro.
    \item Cobertura de UI: agregar tests de comportamiento para pantallas y subcomponentes.
\end{enumerate}

\section{Antes de hacer cualquier cambio}

\begin{enumerate}[leftmargin=*]
    \item Leer la documentacion completa (README.md + docs/).
    \item Leer \texttt{docs/PROJECT\_STATUS.md} para entender que esta activo y que esta diferido.
    \item Leer el proposal del cambio si vas a abrir un nuevo SDD. El flujo esta documentado en \texttt{docs/CONVENTIONS.md} seccion "SDD workflow".
    \item Hablar con el equipo de EPC para alinear prioridades antes de proponer nuevos cambios.
\end{enumerate}

% ============================================================
% Capitulo 11: Convenciones del proyecto
% ============================================================

\chapter{Convenciones del proyecto}

\section{Idioma y estilo}

\begin{itemize}[leftmargin=*]
    \item Espanol rioplatense (voseo) en codigo, comentarios, docs y conversacion.
    \item Ingles solo para mensajes de commit y nombres de variables/funciones publicas de API.
\end{itemize}

\section{Conventional commits}

\begin{itemize}[leftmargin=*]
    \item Mensajes en ingles, formato conventional commits (\texttt{type(scope): subject}).
    \item Subject $\leq$ 50 caracteres.
    \item Body solo cuando el "why" no es obvio.
    \item Sin \texttt{Co-Authored-By}.
    \item Sin emojis en el mensaje.
\end{itemize}

Tipos comunes: \texttt{feat}, \texttt{fix}, \texttt{refactor}, \texttt{docs}, \texttt{test}, \texttt{chore}, \texttt{perf}.

\section{Branches}

\begin{itemize}[leftmargin=*]
    \item \texttt{main}: produccion, consolidada.
    \item \texttt{desarrollo}: integracion continua.
    \item \texttt{merge-desarrollo}: limpieza de PR.
    \item \texttt{feat/<descripcion>}, \texttt{fix/<descripcion>}, \texttt{refactor/<descripcion>}, \texttt{docs/<descripcion>}: feature branches.
\end{itemize}

\section{SDD workflow}

Para cambios sustanciales, seguir el flujo:

\begin{verbatim}
proposal -> specs -> design -> tasks -> apply -> verify -> archive
\end{verbatim}

\begin{enumerate}[leftmargin=*]
    \item \textbf{Proposal} (\texttt{openspec/changes/\{nombre\}/proposal.md}): intent + scope + enfoque.
    \item \textbf{Specs} (\texttt{openspec/changes/\{nombre\}/specs/}): requisitos con Given/When/Then + RFC 2119.
    \item \textbf{Design} (\texttt{openspec/changes/\{nombre\}/design.md}): decisiones B/B/B (Bueno/Barato/Breve) + tradeoffs.
    \item \textbf{Tasks} (\texttt{openspec/changes/\{nombre\}/tasks.md}): checklist numerado completable por sesion.
    \item \textbf{Apply}: implementacion con TDD-strict (cada RED commiteado antes del GREEN).
    \item \textbf{Verify}: code review multi-axis (5 ejes: correctness, readability, architecture, security, performance).
    \item \textbf{Archive}: move a \texttt{openspec/changes/archive/YYYY-MM-DD-\{nombre\}/}.
\end{enumerate}

\textbf{Strict TDD mode} habilitado en \texttt{openspec/config.yaml}. Cada test rojo se commitea antes del feature.

% ============================================================
% Apendices
% ============================================================

\appendix

\chapter{Glosario}

\begin{description}[leftmargin=*]
    \item[EPC] Empresas Publicas de Cundinamarca.
    \item[CRA] Comision de Regulacion de Agua Potable y Saneamiento Basico (Colombia).
    \item[CF] Cargo Fijo. Cobro mensual fijo por la disponibilidad del servicio.
    \item[CC] Cargo por Consumo. Cobro variable segun m\textsuperscript{3} consumidos.
    \item[CMOG] Costo Medio Operativo General.
    \item[CMAA] Costo Medio de Inversiones Ambientales (Res CRA 907/2019).
    \item[IPC] Indice de Precios al Consumidor.
    \item[ISE] Indice de Satisfaccion del Estrato.
    \item[SDD] Spec-Driven Development. Workflow estructurado de desarrollo.
    \item[TDD] Test-Driven Development. Desarrollo guiado por tests.
    \item[RED/GREEN/REFACTOR] Ciclo del TDD.
    \item[B/B/B] Bueno/Barato/Breve. Formato de decision arquitectonica.
    \item[CRUD] Create, Read, Update, Delete.
    \item[FK] Foreign Key.
    \item[DI] Dependency Injection.
    \item[DI] Dependency Injection (inyeccion de dependencias).
\end{description}

\chapter{Referencias normativas}

\begin{enumerate}[leftmargin=*]
    \item \textbf{Res CRA 825/2017.} Por la cual se establece la metodologia tarifaria para las personas prestadoras de los servicios publicos domiciliarios de acueducto y alcantarillado con menos de 5.000 suscriptores.
    \item \textbf{Res CRA 907/2019.} Por la cual se modifican los articulos 9 y 10 de la Res CRA 825/2017. Introduce CMAA y CMVIAA.
    \item \textbf{Res CRA 750/2016.} Por la cual se establece el consumo basico por altitud.
    \item \textbf{Ley 142/1994 art. 99.6.} Ley de Servicios Publicos Domiciliarios. Topes de subsidios y contribuciones.
    \item \textbf{Res CRA 881/2019.} Inversiones ambientales. Adiciona articulos.
\end{enumerate}

\end{document}
